\subsection{EXAMPLES OF FINAL STRUCTURES}

I. Direct image of a structure. When I is a set consisting of a single element, the final structure with respect to \((\mathbf{A}, \mathcal{S}, f)\) is called the direct image under f of the structure G (when it exists).

\begin{enumerate}
\def\labelenumi{\Roman{enumi}.}
\setcounter{enumi}{1}
\tightlist
\item
  Quotient structure. Let A be a set endowed with a structure S of species \(\Sigma\) , let R be an equivalence relation on A, and let \(\varphi\) be the canonical mapping of A onto the quotient set E = A/R (Chapter II, §6, no. 2). The direct image of the structure S under the mapping φ is called (when it exists) the quotient of the structure S by the relation R.
\end{enumerate}

* In general, an order structure or an algebraic structure does not admit quotient structures with respect to arbitrary equivalence relations (cf.~Chapter III, § 1, Exercise 2). On the other hand, a topology always admits a quotient structure with respect to an arbitrary equivalence relation, but this is not necessarily the case for a Hausdorff topology. *

Let A, B be two sets endowed respectively with structures \(\mathcal{S}\) , \(\mathcal{S}'\) of species \(\Sigma\) , and let \(f\) be a morphism of A into B. Let R be the equivalence relation \(f(x) = f(y)\) , let \(\varphi\) be the canonical mapping of A onto A/R, and let \(j\) be the canonical injection of \(f(A)\) into B. Suppose that \(\mathcal{S}\) admits a quotient structure \(\mathcal{S}_0\) with respect to R, and that \(\mathcal{S}'\) induces a structure \(\mathcal{S}_0'\) on \(f(A)\) . Then, in the canonical decomposition \(f = j \circ g \circ \varphi\) of \(f\) (Chapter II, §6, no. 5), the bijection \(g\) of A/R onto \(f(A)\) which is associated with \(f\) is a morphism (but not necessarily an isomorphism) when A/R is endowed with \(\mathcal{S}_0\) and \(f(A)\) with \(\mathcal{S}_0'\) . For \(j \circ g\) is a morphism of A/R into B by the definition of quotient structure, and \(g\) is therefore a morphism of A/R onto \(f(A)\) by the definition of induced structure.

CST20. Let A, A' be two sets endowed with structures \(\mathcal{S}\) , \(\mathcal{S}'\) of species \(\Sigma\) , and let R (resp. R') be an equivalence relation on A (resp. A'). Suppose that there exists a quotient structure \(\mathcal{S}_0\) (resp. \(\mathcal{S}_0'\) ) of \(\mathcal{S}\) by R (resp. \(\mathcal{S}'\) by R'). If \(f\) is a morphism of A into A' which is compatible with the equivalence relations R and R', and if \(g\) is the mapping obtained from \(f\) by passing to the quotients, then \(g\) is a morphism of A/R into A'/R'.

Let \(\varphi\) (resp. \(\varphi'\) ) be the canonical mapping of A onto A/R (resp. of A' onto A'/R'); then we have \(g \circ \varphi = \varphi' \circ f\) . Since \(\varphi'\) and \(f\) are morphisms, so is \(\varphi' \circ f\) by (MO \(_{II}\) ). But then, \(g \circ \varphi\) being a morphism, \(g\) is also a morphism by the definition of quotient structure.

¶ The transitivity criterion CST19 gives rise in particular to the following criterion :

CST21. Let A be a set endowed with a structure \(\mathcal{S}\) of species \(\Sigma\) , and let R be an equivalence relation on A such that there exists on A/R a quotient structure \(\mathcal{S}'\) of \(\mathcal{S}\) by R. Let S be an equivalence relation on A which is coarser than R, and let S/R denote the equivalence relation on A/R which is the quotient of S by R (Chapter II, §6, no. 7). Then there exists on (A/R)/(S/R) a quotient structure \(\mathcal{S}''\) of \(\mathcal{S}'\) by S/R if and only if there exists on A/S a quotient structure \(\mathcal{S}_0\) of \(\mathcal{S}\) by S, and the canonical mapping of A/S (endowed with \(\mathcal{S}_0\) ) onto (A/R)/(S/R) (endowed with \(\mathcal{S}''\) ) is an isomorphism.

Let \(\varphi\) be the canonical mapping of A onto A/R, and let \(\psi\) be that of A/R onto (A/R)/(S/R). By virtue of CST19, \(\mathcal{G}''\) is the quotient of \(\mathcal{G}'\) by S/R if and only if \(\mathcal{G}''\) is the final structure with respect to (A, \(\mathcal{G}\) , \(\psi \circ \varphi\) ). The criterion then follows from the fact that the relation \(\psi(\varphi(x)) = \psi(\varphi(y))\) is equivalent to \(S\{x,y\}\).

Remark. Let A be a set endowed with a structure \(\mathcal{G}\) of species \(\Sigma\) , and let R be an equivalence relation on A such that there exists on E = A/R a quotient structure \(\mathcal{G}'\) of \(\mathcal{G}\) by R. Let \(\varphi\) be the canonical mapping of A onto E. In general, there exists no section \(s\) of \(\varphi\) (Chapter II, §3, no. 8) which is a morphism of E into A. Let us suppose that such a section \(s\) exists, and moreover that there exists a structure \(\mathcal{G}''\) induced by \(\mathcal{G}\) on \(s(E)\) . Then, if \(j\) denotes the canonical injection of \(s(E)\) into A and if \(s = j \circ f\) , the bijection \(f\) is an isomorphism of E onto \(s(E)\) . For \(f\) is a morphism by the definition of induced structure, and \(g = \varphi \circ j\) is a morphism of \(s(E)\) onto E by reason of \((MO_{II})\) . Since \(g \circ f\) and \(f \circ g\) are the identity mappings of E and \(s(E)\) , respectively, the assertion is a consequence of CST8.

